\section{논의 및 한계}
\label{sec:discussion}

\noindent\textbf{관측 확장의 효과와 비용.}
다해상도 결합은 총 토큰 상한을 맞춘 단일 해상도 대조군보다 높은 \AUPRO를 보였고, 관측별 독립 매칭도 특징 연결보다 높은 성능을 보였다.
반면 여러 해상도를 추론하므로 단일 해상도보다 쿼리 시간과 메모리가 늘어난다.

\noindent\textbf{보정과 집계의 효과 범위.}
Q95 보정과 Rényi 집계의 결합 이득은 \AUPRO에서 나타났으며 \pAUROC로 이어지지 않았다.
클래스별로 효과의 방향이 다르고 fabric의 기여가 크므로, 이 이득을 모든 결함 유형에 대한 일관된 개선으로 볼 수는 없다.
또한 정상 영역의 추가 응답 감소는 점수 분포의 변화이며, 고정 임계값에서의 오경보율 감소를 직접 보여주지는 않는다.

\noindent\textbf{서포트 편차와 배포 임계값.}
관측 확장은 주어진 정상 표본의 비교 단서를 늘리지만 미관측 정상 외관을 직접 보지는 못하므로, 서포트 선택에 따른 성능 변동이 남는다(1-shot 표준편차 4.07\,pp).
또한 \AUPRO와 \pAUROC는 FPR 구간에 대한 적분 지표이므로, 실제 라인에 적용하려면 목표 오경보율에 맞춘 임계값 교정이 별도로 필요하다.
